% GENERATED FILE. Edit canonical lessons, then run npm run generate:books.
% canonical-source-hash: fnv1a64:1ba04cc1de171f38
% canonical-lessons: PA-C56-mahiman, PA-C56-pina, PA-C56-sunna, PA-C56-baithna, PA-C56-sauna

\chapter{Guest, To drink, To hear, To sit, To sleep}
\label{ch:pa-guest-to-drink-to-hear-to-sit-to-sleep}
\hlchaptermodality{56}

\begin{chapteropening}
\textbf{By the end of this chapter:} \emph{I can say a guest, to drink, to hear, to listen, to sit and to sleep.}

Complete the last lesson of chapter 56: i can say a guest, to drink, to hear, to listen, to sit and to sleep.
\end{chapteropening}

\section[mahimān]{\pa{ਮਹਿਮਾਨ} (mahimān) — a guest}

\label{lesson:PA-C56-mahiman}



\begin{quote}
Before the new one: say the Punjabi for a girl, then the Punjabi for a man.
\end{quote}

\subsection*{You'll want to know: \pa{ਮਹਿਮਾਨ}}
\textbf{\pa{ਮਹਿਮਾਨ}} — \emph{mahimān} — \textquotedblleft{}a guest\textquotedblright{}.

\begin{grammarlens}[title={What you've built}]
One more word for this chapter.
\end{grammarlens}

\subsection*{Your turn}
\begin{itemize}
\raggedright
  \item \emph{Say it:} \emph{mahimān}
  \item \emph{Say it:} \emph{mahimān}, once more
  \item \emph{Say it:} say \emph{mahimān}
  \item \emph{Recall:} say the Punjabi for a girl, then the Punjabi for a man, then say \emph{mahimān} again
\end{itemize}

\begin{tcolorbox}[breakable,colback=teal!4,colframe=teal!35!black,title={Before you move on}]
What does \pa{ਮਹਿਮਾਨ} mean? (\textquotedblleft{}A guest\textquotedblright{}.) Say it once more.
\end{tcolorbox}
\section[pīṇā]{\pa{ਪੀਣਾ} (pīṇā) — to drink}

\label{lesson:PA-C56-pina}



\begin{quote}
Before the new one: say the Punjabi for a man, then the Punjabi for a guest.
\end{quote}

\subsection*{You'll want to know: \pa{ਪੀਣਾ}}
\textbf{\pa{ਪੀਣਾ}} — \emph{pīṇā} — \textquotedblleft{}to drink\textquotedblright{}.

\begin{grammarlens}[title={What you've built}]
One more word for this chapter.
\end{grammarlens}

\subsection*{Your turn}
\begin{itemize}
\raggedright
  \item \emph{Say it:} \emph{pīṇā}
  \item \emph{Say it:} \emph{pīṇā}, once more
  \item \emph{Say it:} say \emph{pīṇā}
  \item \emph{Recall:} say the Punjabi for a man, then the Punjabi for a guest, then say \emph{pīṇā} again
\end{itemize}

\begin{tcolorbox}[breakable,colback=teal!4,colframe=teal!35!black,title={Before you move on}]
What does \pa{ਪੀਣਾ} mean? (\textquotedblleft{}To drink\textquotedblright{}.) Say it once more.
\end{tcolorbox}
\section[suṇnā]{\pa{ਸੁਣਨਾ} (suṇnā) — to hear, to listen}

\label{lesson:PA-C56-sunna}



\begin{quote}
Before the new one: say the Punjabi for a guest, then the Punjabi for to drink.
\end{quote}

\subsection*{You'll want to know: \pa{ਸੁਣਨਾ}}
\textbf{\pa{ਸੁਣਨਾ}} — \emph{suṇnā} — \textquotedblleft{}to hear, to listen\textquotedblright{}.

\begin{grammarlens}[title={What you've built}]
One more word for this chapter.
\end{grammarlens}

\subsection*{Your turn}
\begin{itemize}
\raggedright
  \item \emph{Say it:} \emph{suṇnā}
  \item \emph{Say it:} \emph{suṇnā}, once more
  \item \emph{Say it:} say \emph{suṇnā}
  \item \emph{Recall:} say the Punjabi for a guest, then the Punjabi for to drink, then say \emph{suṇnā} again
\end{itemize}

\begin{tcolorbox}[breakable,colback=teal!4,colframe=teal!35!black,title={Before you move on}]
What does \pa{ਸੁਣਨਾ} mean? (\textquotedblleft{}To hear, to listen\textquotedblright{}.) Say it once more.
\end{tcolorbox}
\section[baiṭhṇā]{\pa{ਬੈਠਣਾ} (baiṭhṇā) — to sit}

\label{lesson:PA-C56-baithna}



\begin{quote}
Before the new one: say the Punjabi for to drink, then the Punjabi for to hear.
\end{quote}

\subsection*{You'll want to know: \pa{ਬੈਠਣਾ}}
\textbf{\pa{ਬੈਠਣਾ}} — \emph{baiṭhṇā} — \textquotedblleft{}to sit\textquotedblright{}.

\begin{grammarlens}[title={What you've built}]
One more word for this chapter.
\end{grammarlens}

\subsection*{Your turn}
\begin{itemize}
\raggedright
  \item \emph{Say it:} \emph{baiṭhṇā}
  \item \emph{Say it:} \emph{baiṭhṇā}, once more
  \item \emph{Say it:} say \emph{baiṭhṇā}
  \item \emph{Recall:} say the Punjabi for to drink, then the Punjabi for to hear, then say \emph{baiṭhṇā} again
\end{itemize}

\begin{tcolorbox}[breakable,colback=teal!4,colframe=teal!35!black,title={Before you move on}]
What does \pa{ਬੈਠਣਾ} mean? (\textquotedblleft{}To sit\textquotedblright{}.) Say it once more.
\end{tcolorbox}
\section[sauṇā]{\pa{ਸੌਣਾ} (sauṇā) — to sleep}

\label{lesson:PA-C56-sauna}



\begin{quote}
Before the new one: say the Punjabi for to hear, then the Punjabi for to sit.
\end{quote}

\subsection*{You'll want to know: \pa{ਸੌਣਾ}}
\textbf{\pa{ਸੌਣਾ}} — \emph{sauṇā} — \textquotedblleft{}to sleep\textquotedblright{}.

\begin{grammarlens}[title={What you've built}]
One more word for this chapter.
\end{grammarlens}

\subsection*{Your turn}
\begin{itemize}
\raggedright
  \item \emph{Say it:} \emph{sauṇā}
  \item \emph{Say it:} \emph{sauṇā}, once more
  \item \emph{Say it:} say \emph{sauṇā}
  \item \emph{Recall:} say the Punjabi for to hear, then the Punjabi for to sit, then say \emph{sauṇā} again
\end{itemize}

\begin{tcolorbox}[breakable,colback=teal!4,colframe=teal!35!black,title={Before you move on}]
What does \pa{ਸੌਣਾ} mean? (\textquotedblleft{}To sleep\textquotedblright{}.) Say it once more.
\end{tcolorbox}
